\section{Beyond the Torah: A Corpus Accumulates}

No ancient source records a second commission. The rest of the Greek Old
Testament came into being the way large literatures usually do: book by
book, over roughly two centuries, by different hands, in more than one
place, without a controlling authority---and the one contemporary witness we
possess says so plainly.

That witness is the prologue that the grandson of Jesus ben Sira attached
to his Greek translation of his grandfather's Hebrew wisdom book
(Ecclesiasticus/Sirach). Writing after his arrival in Egypt in the
thirty-eighth year of King Euergetes---on the usual reckoning,
132~B.C.---he apologizes for his own rendering and observes that ``the law
itself, and the prophets, and the rest of the books, have no small
difference, when they are spoken in their own language,'' since things
first uttered in Hebrew ``have not the same force in them'' when
translated.\endnote{Prologue to Sirach, quoted from Brenton's public-domain
English (1851), checked at ebible.org. The prologue's threefold
``law\ldots prophets\ldots rest of the books'' is also a key early witness
to the shape of the Hebrew collection itself.} Four facts stand in this
single paragraph: by the later second century B.C. the prophets and other
books already existed in Greek; individuals (here a named one) made such
translations on their own initiative; contemporaries knew the versions
differed in quality from the originals; and ``the Septuagint'' as a bounded,
authorized whole is nowhere in view.

The corpus that eventually sailed under the name confirms this piecemeal
origin by its internal variety. Rendering styles run from the studied
literalism that tracks Hebrew particles into Greek, to the free literary
manner of Job and Proverbs, to outright rewriting; some books circulated in
two Greek forms (Daniel, Judges); some Greek books are expansions of their
Hebrew base (the additions to Esther and Daniel); and some---the Wisdom of
Solomon, 2~Maccabees---were composed in Greek and translate
nothing.\endnote{Book-level characterizations follow the standard modern
surveys and edition introductions listed in the References (class~K); the
argument here needs only the uncontested variety, not any single book's
profile.} ``Septuagint,'' extended from the Torah legend to this whole
shelf, is thus a Christian-era convenience label. Scholars mark the
distinction by reserving ``Old Greek'' for the earliest recoverable
translation of a given book---an editorial construct, not a surviving
manuscript---and this account uses the name ``Septuagint'' for the corpus
only in that loose, received sense.
